
\section{Appendix B. Implementation Details}

\subsection{B.1 The SVDLinear Module}

\begin{lstlisting}
class SVDLinear(nn.Module):
    def \_\_init\_\_(self, in\_features, out\_features, rank, lora\_rank=16, ...):
        super().\_\_init\_\_()
        self.V = nn.Parameter(torch.empty(rank, in\_features))
        self.U = nn.Parameter(torch.empty(out\_features, rank))
        if bias:
            self.bias = nn.Parameter(torch.empty(out\_features))
        if lora\_rank \textgreater{} 0:
            self.lora\_A = nn.Parameter(torch.empty(lora\_rank, in\_features))
            self.lora\_B = nn.Parameter(torch.zeros(out\_features, lora\_rank))
    
    def forward(self, x):
        Vx = F.linear(x, self.V)
        out = F.linear(Vx, self.U, self.bias)
        if self.lora\_rank \textgreater{} 0:
            out = out + F.linear(F.linear(x, self.lora\_A), self.lora\_B) * self.scaling
        return out
\end{lstlisting}

\subsection{B.2 The Decay Gate}

\begin{lstlisting}
def \_init\_decay\_gate(self):
    with torch.no\_grad():
        nn.init.kaiming\_uniform\_(self.decay\_gate.weight, a=math.sqrt(5))
        if self.decay\_channel\_wise:
            decay\_times = torch.exp(torch.linspace(math.log(1.0), math.log(1024.0), self.num\_heads))
            decay\_times = decay\_times.repeat\_interleave(self.head\_dim)
        else:
            decay\_times = torch.exp(torch.linspace(math.log(1.0), math.log(1024.0), self.num\_heads))
        bias\_target = -torch.log(decay\_times)
        self.decay\_gate.bias.copy\_(bias\_target)
\end{lstlisting}

\textbf{Note.} As shown in Appendix A.6, this initialization is \textbf{inverted}. 
The intended $\tau_h$ values are $[1, 1024]$, but the actual $\tau_h = 1 + 1/\tau_h^\text{target}$ 
ranges from 1.001 to 2. Heads converge during training.

\subsection{B.3 The Chunked Scan}

\begin{lstlisting}
def \_forward\_decay\_gated\_linear\_kernel(self, q, k, v, hidden\_states, mask=None, past\_state=None):
    B, H, N, D\_k = q.shape
    D\_v = v.shape[-1]
    
    phi\_q = F.elu(q.float()) + 1.0
    phi\_k = F.elu(k.float()) + 1.0
    v\_f = v.float()
    
    if mask is not None:
        phi\_k = phi\_k * mask\_expanded
    
    # Decay
    decay\_logits = self.decay\_gate(hidden\_states)
    if self.decay\_channel\_wise:
        decay\_logits = decay\_logits.view(B, N, H, D\_k).permute(0, 2, 1, 3)
        lambda\_decay = compute\_exponential\_decay(decay\_logits).unsqueeze(-1)
    else:
        decay\_logits = decay\_logits.permute(0, 2, 1).unsqueeze(-1).unsqueeze(-1)
        lambda\_decay = compute\_exponential\_decay(decay\_logits)
    
    # Force padding to have $\lambda$=1
    if mask is not None:
        lambda\_decay = torch.where(mask\_decay \textgreater{} 0, lambda\_decay, torch.ones\_like(lambda\_decay))
    
    # Chunked scan
    chunk\_size = 64
    num\_chunks = (N + chunk\_size - 1) // chunk\_size
    S\_chunk = torch.zeros(B, H, D\_k, D\_v, dtype=torch.float32)
    Z\_chunk = torch.zeros(B, H, D\_k, 1, dtype=torch.float32)
    outputs = []
    
    for c in range(num\_chunks):
        s, e = c * chunk\_size, min(N, (c+1) * chunk\_size)
        qc, kc, vc = phi\_q[:,:,s:e], phi\_k[:,:,s:e], v\_f[:,:,s:e]
        lam\_c = lambda\_decay[:,:,s:e]
        Lc = e - s
        
        # Intra-chunk
        log\_lam = torch.log(lam\_c.clamp(min=1e-6))
        log\_gamma = torch.cumsum(log\_lam, dim=-1)
        log\_M = log\_gamma.unsqueeze(-1) - log\_gamma.unsqueeze(-2)
        M\_intra = torch.exp(log\_M.clamp(max=0.0)) * causal\_tril
        aff = torch.matmul(qc, kc.transpose(-1, -2)) * M\_intra
        num\_intra = torch.matmul(aff, vc)
        den\_intra = aff.sum(dim=-1, keepdim=True)
        
        # Inter-chunk
        decay = torch.exp(log\_gamma.clamp(max=0.0)).unsqueeze(-1)
        num\_inter = torch.matmul(qc, S\_chunk) * decay
        den\_inter = torch.matmul(qc, Z\_chunk) * decay
        
        out\_c = (num\_intra + num\_inter) / (den\_intra + den\_inter + self.eps)
        outputs.append(out\_c)
        
        # Update state
        total\_decay = torch.exp(log\_gamma[:,:,-1:].clamp(max=0.0)).unsqueeze(-1)
        w\_to\_end = torch.exp((log\_gamma[:,:,-1:] - log\_gamma).clamp(max=0.0)).unsqueeze(-1)
        kc\_w = kc * w\_to\_end
        S\_chunk = total\_decay * S\_chunk + torch.matmul(kc\_w.transpose(-1, -2), vc)
        Z\_chunk = total\_decay * Z\_chunk + kc\_w.sum(dim=-2, keepdim=True).transpose(-1, -2)
    
    context = torch.cat(outputs, dim=2).to(q.dtype)
    return context, (S\_chunk, Z\_chunk)
\end{lstlisting}

\subsection{B.4 The `step\_recurrent` Method}

\begin{lstlisting}
def step\_recurrent(self, x\_t, past\_state=None):
    if x\_t.dim() == 2:
        x\_t = x\_t.unsqueeze(1)
    
    B = x\_t.shape[0]
    q = self.q\_proj(x\_t).view(B, 1, self.num\_heads, self.head\_dim).transpose(1, 2)
    k = self.k\_proj(x\_t).view(B, 1, self.num\_kv\_heads, self.head\_dim).transpose(1, 2)
    v = self.v\_proj(x\_t).view(B, 1, self.num\_kv\_heads, self.head\_dim).transpose(1, 2)
    
    if self.num\_kv\_heads \textless{} self.num\_heads:
        k = repeat\_kv(k, self.kv\_repeat)
        v = repeat\_kv(v, self.kv\_repeat)
    
    phi\_q = F.elu(q.float()) + 1.0
    phi\_k = F.elu(k.float()) + 1.0
    
    # Decay
    if self.use\_decay\_gating:
        decay\_logits = self.decay\_gate(x\_t)
        # ... compute $\lambda$\_t ...
    else:
        lam\_t = torch.ones(B, H, 1, 1)
    
    # Update state
    if past\_state is not None:
        S\_prev, Z\_prev = past\_state
    else:
        S\_prev = torch.zeros(B, H, D\_k, D\_v)
        Z\_prev = torch.zeros(B, H, D\_k, 1)
    
    kt\_t = phi\_k.transpose(-1, -2)
    kv\_t = torch.matmul(kt\_t, v.float())
    
    S\_t = lam\_t * S\_prev + kv\_t
    Z\_t = lam\_t * Z\_prev + kt\_t
    
    num\_t = torch.matmul(phi\_q, S\_t)
    den\_t = torch.matmul(phi\_t, Z\_t) + self.eps
    y\_t = num\_t / den\_t
    
    # Output projection
    context = y\_t.transpose(1, 2).contiguous().view(B, 1, self.all\_head\_size)
    if self.is\_nested\_encoder:
        out\_t = context
    else:
        out\_t = self.out\_proj(context)
    
    return out\_t, (S\_t, Z\_t)
\end{lstlisting}

\subsection{B.5 Common Pitfalls}

\subsubsection{B.5.1 Decay Init Is Inverted}

As noted in A.6, the decay bias initialization is inverted. Consider 
fixing to \texttt{bias\_target = -softplus\_inv(log(decay\_times))}.

\subsubsection{B.5.2 Padding in Feature Map}

 (x) = ELU(x) + 1 \textgreater{} 0 always, so padding tokens contribute to the state. 
We mask them with \texttt{phi\_k = phi\_k * mask\_expanded}.

\subsubsection{B.5.3 Chunk Size}

Fixed at 64. Larger chunks give more parallelism but reduce decay effect.

\subsubsection{B.5.4 Nested Encoder Output Projection}

For RoBERTa, \texttt{is\_nested\_encoder = True} means the output projection is 
skipped (applied externally).
\begin{lstlisting}

---
\end{lstlisting}